\subsection{retee}
{
The final example is a program called “re-tee.” It takes three arguments: the pathname of an executable program, the
process id of a running instance of the same program, and a file name. It adds code to the running program that
copies to the named file all output that the program writes to its standard output file descriptor. In this way it
works like “tee,” which passes output along to its own standard out while also saving it in a file. The motivation
for the example program is that you run a program, and it starts to print copious lines of output to your screen, and
you wish to save that output in a file without having to re-run the program.


\bigskip

{
\#include <stdio.h>

{
\#include <fcntl.h>

{
\#include <vector>

{
\#include {\textquotedbl}BPatch.h{\textquotedbl}

{
\#include {\textquotedbl}BPatch_point.h{\textquotedbl}}

{
\#include {\textquotedbl}BPatch_process.h{\textquotedbl}}

{
\#include {\textquotedbl}BPatch_function.h{\textquotedbl}}

{
\#include {\textquotedbl}BPatch_thread.h{\textquotedbl}


\bigskip

{
/*

{
* retee.C}

{
*}

{
* This program (mutator) provides an example of several facets of}

{
* Dyninst's behavior, and is a good basis for many Dyninst}

{
* mutators. We want to intercept all output from a target application}

{
* (the mutatee), duplicating output to a file as well as the}

{
* original destination (e.g., stdout).}

{
*}

{
* This mutator operates in several phases. In brief:}

{
* 1) Attach to the running process and get a handle (BPatch_process}

{
* object)}

{
* 2) Get a handle for the parsed image of the mutatee for function}

{
* lookup (BPatch_image object)}

{
* 3) Open a file for output}

{
* 3a) Look up the {\textquotedbl}open{\textquotedbl} function}

{
* 3b) Build a code snippet to call open with the file name.}

{
* 3c) Run that code snippet via a oneTimeCode, saving the returned}

{
* file descriptor}

{
* 4) Write the returned file descriptor into a memory variable for}

{
* mutatee-side use}

{
* 5) Build a snippet that copies output to the file}

{
* 5a) Locate the {\textquotedbl}write{\textquotedbl} library call}

{
* 5b) Access its parameters}

{
* 5c) Build a snippet calling write(fd, parameters)}

{
* 5d) Insert the snippet at write}

{
* 6) Add a hook to exit to ensure that we close the file (using}

{
* a callback at exit and another oneTimeCode)}

{
*/}


\bigskip

{
void usage() {}

{
fprintf(stderr, {\textquotedbl}Usage: retee <process pid>
<filename>{\textbackslash}n{\textquotedbl});}

{
fprintf(stderr, {\textquotedbl} note: <filename> is relative to the application
process.{\textbackslash}n{\textquotedbl});}

{
}}


\bigskip

{
// We need to use a callback, and so the things that callback requires}

{
// are made global - this includes the file descriptor snippet (see below)}

{
BPatch_variableExpr *fdVar = NULL;}


\bigskip

{
// Before we add instrumentation, we need to open the file for}

{
// writing. We can do this with a oneTimeCode - a piece of code run at}

{
// a particular time, rather than at a particular location.}


\bigskip

{
int openFileForWrite(BPatch_process *app, BPatch_image *appImage, char *fileName) {}

{
// The code to be generated is:

{
// fd = open(argv[2], O_WRONLY{\textbar}O_CREAT, 0666);


\bigskip

{
// (1) Find the open function

{
std::vector<BPatch_function *>openFuncs;

{
appImage->findFunction({\textquotedbl}open{\textquotedbl}, openFuncs);

{
if (openFuncs.size() == 0) {}

{
fprintf(stderr, {\textquotedbl}ERROR: Unable to find function for
open(){\textbackslash}n{\textquotedbl});}

{
return -1;}

{
}}


\bigskip

{
// (2) Allocate a vector of snippets for the parameters to open

{
std::vector<BPatch_snippet *> openArgs;


\bigskip

{
// (3) Create a string constant expression from argv[3]

{
BPatch_constExpr fileNameExpr(fileName);


\bigskip

{
// (4) Create two more constant expressions _WRONLY{\textbar}O_CREAT and 0666

{
BPatch_constExpr fileFlagsExpr(O_WRONLY{\textbar}O_CREAT);

{
BPatch_constExpr fileModeExpr(0666);


\bigskip

{
// (5) Push 3 & 4 onto the list from step 2, push first to last parameter.}

{
openArgs.push_back(&fileNameExpr);

{
openArgs.push_back(&fileFlagsExpr);

{
openArgs.push_back(&fileModeExpr);


\bigskip

{
// (6) create a procedure call using function found at 1 and

{
// parameters from step 5.

{
BPatch_funcCallExpr openCall(*openFuncs[0], openArgs);


\bigskip

{
// (7) The oneTimeCode returns whatever the return result from}

{
// the BPatch_snippet is. In this case, the return result of}

{
// open -> the file descriptor.}

{
void *openFD = app->oneTimeCode( openCall );}


\bigskip

{
// oneTimeCode returns a void *, and we want an int file handle}

{
return (int) (long) openFD;}

{
}}


\bigskip

{
// We have used a oneTimeCode to open the file descriptor. However,}

{
// this returns the file descriptor to the mutator - the mutatee has}

{
// no idea what the descriptor is. We need to allocate a variable in}

{
// the mutatee to hold this value for future use and copy the}

{
// (mutator-side) value into the mutatee variable.}


\bigskip

{
// Note: there are alternatives to this technique. We could have}

{
// allocated the variable before the oneTimeCode and augmented the}

{
// snippet to do the assignment. We could also write the file}

{
// descriptor as a constant into any inserted instrumentation.}


\bigskip

{
BPatch_variableExpr *writeFileDescIntoMutatee(BPatch_process *app,

{
BPatch_image *appImage,

{
int fileDescriptor) {}

{
// (1) Allocate a variable in the mutatee of size (and type) int}

{
BPatch_variableExpr *fdVar =
app->malloc(*appImage->findType({\textquotedbl}int{\textquotedbl}));

{
if (fdVar == NULL) return NULL;}


\bigskip

{
// (2) Write the value into the variable}

{
// Like memcpy, writeValue takes a pointer}

{
// The third parameter is for functionality called {\textquotedbl}saveTheWorld{\textquotedbl},

{
// which we don't worry about here (and so is false)}

{
bool ret = fdVar->writeValue((void *) &fileDescriptor, sizeof(int),}

{
 false);

{
if (ret == false) return NULL;}


\bigskip

{
return fdVar;}

{
}}


\bigskip

{
// We now have an open file descriptor in the mutatee. We want to}

{
// instrument write to intercept and copy the output. That happens}

{
// here.


\bigskip

{
bool interceptAndCloneWrite(BPatch_process *app,

{
BPatch_image *appImage,

{
BPatch_variableExpr *fdVar) {}

{
// (1) Locate the write call}

{
std::vector<BPatch_function *> writeFuncs;}


\bigskip

{
appImage->findFunction({\textquotedbl}write{\textquotedbl},

{
writeFuncs);}

{
if(writeFuncs.size() == 0) {}

{
fprintf(stderr, {\textquotedbl}ERROR: Unable to find function for
write(){\textbackslash}n{\textquotedbl});}

{
return false;}

{
}}


\bigskip

{
// (2) Build the call to (our) write. Arguments are:}

{
// ours: fdVar (file descriptor)}

{
// parameter: buffer}

{
// parameter: buffer size}


\bigskip

{
// Declare a vector to hold these.}

{
std::vector<BPatch_snippet *> writeArgs;}

{
// Push on the file descriptor}

{
writeArgs.push_back(fdVar);}

{
// Well, we need the buffer... but that's a parameter to the}

{
// function we're implementing. That's not a problem - we can grab}

{
// it out with a BPatch_paramExpr.}

{
BPatch_paramExpr buffer(1); // Second (0, 1, 2) argument}

{
BPatch_paramExpr bufferSize(2);}

{
writeArgs.push_back(&buffer);}

{
writeArgs.push_back(&bufferSize);}


\bigskip

{
// And build the write call}

{
BPatch_funcCallExpr writeCall(*writeFuncs[0], writeArgs);}


\bigskip

{
// (3) Identify the BPatch_point for the entry of write. We're}

{
// instrumenting the function with itself; normally the findPoint}

{
// call would operate off a different function than the snippet.}


\bigskip

{
std::vector<BPatch_point *> *points;}

{
points = writeFuncs[0]->findPoint(BPatch_entry);}

{
if ((*points).size() == 0) {

{
return false;}

{
}


\bigskip

{
// (4) Insert the snippet at the start of write}


\bigskip

{
return app->insertSnippet(writeCall, *points);}


\bigskip

{
// Note: we have just instrumented write() with a call to}

{
// write(). This would ordinarily be a _bad thing_, as there is}

{
// nothing to stop infinite recursion - write -> instrumentation}

{
// -> write -> instrumentation....}

{
// However, Dyninst uses a feature called a {\textquotedbl}tramp guard{\textquotedbl} to}

{
// prevent this, and it's on by default.

{
}}


\bigskip

{
// This function is called as an exit callback (that is, called}

{
// immediately before the process exits when we can still affect it)}

{
// and thus must match the exit callback signature:}

{
//

{
// typedef void (*BPatchExitCallback) (BPatch_thread *, BPatch_exitType)}

{
//

{
// Note that the callback gives us a thread, and we want a process - but}

{
// each thread has an up pointer.}


\bigskip

{
void closeFile(BPatch_thread *thread, BPatch_exitType) {}

{
fprintf(stderr, {\textquotedbl}Exit callback called for process...{\textbackslash}n{\textquotedbl});}


\bigskip

{
// (1) Get the BPatch_process and BPatch_images}

{
BPatch_process *app = thread->getProcess();}

{
BPatch_image *appImage = app->getImage();}


\bigskip

{
// The code to be generated is:}

{
// close(fd);}


\bigskip

{
// (2) Find close}

{
std::vector<BPatch_function *> closeFuncs;

{
appImage->findFunction({\textquotedbl}close{\textquotedbl}, closeFuncs);

{
if (closeFuncs.size() == 0) {}

{
fprintf(stderr, {\textquotedbl}ERROR: Unable to find function for
close(){\textbackslash}n{\textquotedbl});}

{
return;}

{
}}


\bigskip

{
// (3) Allocate a vector of snippets for the parameters to open

{
std::vector<BPatch_snippet *> closeArgs;


\bigskip

{
// (4) Add the fd snippet - fdVar is global since we can't}

{
// get it via the callback}

{
closeArgs.push_back(fdVar);}


\bigskip

{
// (5) create a procedure call using function found at 1 and

{
// parameters from step 3.

{
BPatch_funcCallExpr closeCall(*closeFuncs[0], closeArgs);


\bigskip

{
// (6) Use a oneTimeCode to close the file}

{
app->oneTimeCode( closeCall );}


\bigskip

{
// (7) Tell the app to continue to finish it off.}

{
app->continueExecution();}


\bigskip

{
return;}

{
}}


\bigskip

{
BPatch bpatch;


\bigskip

{
// In main we perform the following operations.}

{
// 1) Attach to the process and get BPatch_process and BPatch_image}

{
// handles}

{
// 2) Open a file descriptor}

{
// 3) Instrument write}

{
// 4) Continue the process and wait for it to terminate}


\bigskip

{
int main(int argc, char *argv[]) {

{
int pid;

{
if (argc != 3) {

{
usage();}

{
exit(1);

{
}

{
pid = atoi(argv[1]);


\bigskip

{
// Attach to the program - we can attach with just a pid; the}

{
// program name is no longer necessary}

{
fprintf(stderr, {\textquotedbl}Attaching to process \%d...{\textbackslash}n{\textquotedbl}, pid);}

{
BPatch_process *app = bpatch.processAttach(NULL, pid);


\bigskip

{
if (!app) return -1;}


\bigskip

{
// Read the program's image and get an associated image object

{
BPatch_image *appImage = app->getImage();

{
std::vector<BPatch_function*> writeFuncs;


\bigskip

{
fprintf(stderr, {\textquotedbl}Opening file \%s for write...{\textbackslash}n{\textquotedbl}, argv[2]);}

{
int fileDescriptor = openFileForWrite(app, appImage, argv[2]);}


\bigskip

{
if (fileDescriptor == -1) {}

{
fprintf(stderr, {\textquotedbl}ERROR: opening file \%s for write
failed{\textbackslash}n{\textquotedbl},}

{
argv[2]);}

{
exit(1);}

{
}}


\bigskip

{
fprintf(stderr, {\textquotedbl}Writing returned file descriptor \%d into{\textquotedbl}}

{
{\textquotedbl}mutatee...{\textbackslash}n{\textquotedbl}, fileDescriptor);}


\bigskip

{
// This was defined globally as the exit callback needs it.}

{
fdVar = writeFileDescIntoMutatee(app, appImage, fileDescriptor);}

{
if (fdVar == NULL) {}

{
fprintf(stderr, {\textquotedbl}ERROR: failed to write mutatee-side
variable{\textbackslash}n{\textquotedbl});}

{
exit(1);}

{
}}


\bigskip

{
fprintf(stderr, {\textquotedbl}Instrumenting write...{\textbackslash}n{\textquotedbl});}

{
bool ret = interceptAndCloneWrite(app, appImage, fdVar);}

{
if (!ret) {}

{
fprintf(stderr, {\textquotedbl}ERROR: failed to instrument mutatee{\textbackslash}n{\textquotedbl});}

{
exit(1);}

{
}}


\bigskip

{
fprintf(stderr, {\textquotedbl}Adding exit callback...{\textbackslash}n{\textquotedbl});}

{
bpatch.registerExitCallback(closeFile);}


\bigskip

{
// Continue the execution...}

{
fprintf(stderr, {\textquotedbl}Continuing execution and waiting for
termination{\textbackslash}n{\textquotedbl});}

{
app->continueExecution();}


\bigskip

{
while (!app->isTerminated())

{
bpatch.waitForStatusChange();}


\bigskip

{
printf({\textquotedbl}Done.{\textbackslash}n{\textquotedbl});}


\bigskip

{
return 0;}

{
}}
